% \section{3D Egocentric Foundation Model (\EFMmodel{})}
\section{Egocentric Voxel Lifting (\EFMmodel{})}

We design \EFMmodel{} as a universal 3D backbone that lifts 2D foundation features from posed and calibrated video streams into a 3D gravity-aligned voxel grid of features. Before processing this 3D feature volume with a 3D U-Net we concatenate masks derived from semi-dense points that indicate surface points as well as freespace. As shown in Fig.~\ref{fig:efm_model}, on this feature volume we can run different kinds of 3D-CNN heads to regress desired 3D quantities such as 3D OBB parameters and occupancy values.
%We propose the first 3D egocentric foundation model,\EFMmodel{}, inspired by related work~\cite{murez2020atlas,philion2020lift,rukhovich2022imvoxelnet,sun2021neucon} to lift 2D foundation features~\cite{oquab2023dinov2,darcet2023vitneedreg} extracted from posed video snippets into an explicit gravity-aligned feature volume. This feature volume is defined on a voxel grid in a local coordinate system. 

\begin{figure}
    \centering
    \includegraphics[width=\textwidth]{figures/efm3d_arch_v1.pdf}
    \caption{\EFMmodel{} lifts 2D features extracted from frozen image foundation models into an local gravity-aligned 3D voxel grid of features. After concatenating point masks a U-Net processes the 3D feature volume before running 3D CNN heads. These heads regress target parameters such as 3D OBBs and occupancy values.
    %\TODO{show gravity aligned voxel grid}
    \label{fig:efm_model}}
\end{figure}

%\subsection{2D-3D Lifting Backbone}
%The 2D-3D lifting backbone is the core \EFMmodel{} module that enables 3D reasoning from 2D posed image observations. It first encodes the images, then lifts those features into 3D. In 3D we concatenate with a semidense points mask before running a 3D U-Net to produce the final 3D feature volume.

\noindent\textbf{2D-3D Lifting Encoder} We run a 2D foundation model~\cite{oquab2023dinov2} with frozen weights on $T$ posed input images from each video stream. The resulting 2D foundation features are up-sampled to the full input image resolution using stacked 2D CNN \& upsample layers.  
%
We instantiate a local, gravity-aligned 3D voxel grid in front of the most recent Aria camera rig. We choose the last RGB camera pose and gravity align its rotation to obtain the anchor pose of the voxel grid.
%\TODO{could go to supplement}
%Given camera rotation $R_{wc} = [r_x, r_y, r_z]$ and unit-length gravity direction $g_w$ in world coordinates we compute the gravity aligned rotation $R_{wg}$ as:
%\begin{align}
%    R_{wg} &= [g_w, \lceil d_z \times g_w \rceil, \lceil d_z\rceil ] \quad \text{where}\quad
%     d_z %= \text{reject}(R_{wc}_2, g_w) 
%     = r_z - g_w^T r_z g_w \,,
%\end{align}
%and $\lceil x \rceil $ normalizes a vector $x$ to unit length.
%
Using known camera poses and calibrations, we project the voxel grid centers into each image and sample the 2D feature maps using bi-linear interpolation while respecting the valid radius of the distorted camera models.
This yields a feature volume of $T\times F \times D\times H\times W$ for every stream. 
We aggregate features across all streams and the time dimension using mean and standard deviation to obtain a lifted feature volume of shape $2F \times D\times H\times W$.
The metric size and resolution of the voxel grid can be defined freely. 
%
Note that this lifting formulation is camera model agnostic and supports online-calibrated video streams from Aria. Please refer to the supplemental for the experiments of \EFMmodel{} using different camera models.

\noindent\textbf{Semi-Dense Point Encoder}. We additionally input the semi-dense 3D point cloud provided by the Aria research toolkit~\cite{engel2023project} to provide a geometric prior for the lifting backbone. Together with the visibility information from where each point was observed this point cloud contains not only information about the surfaces in the scene but also the freespace. To provide this information to \EFMmodel{} we compute a point and a free-space mask tensor to concatenate to the lifted feature volume. 
%After transforming the 3D point positions into the voxel grid coordinate system, 
The point mask indicates where there are surface points. It is computed by discretizing points into a binary mask sized $D\times H\times W$ corresponding to the voxel grid resolution.
The freespace mask indicates the voxel locations between the camera centers to the known surface points.
We compute it by sampling $S$ points along each ray within the bounds of the feature volume. 
We concatenate both masks to the lifted 3D features resulting in a 3D feature volume shaped $(F+2)\times D\times H\times W$.

\noindent\textbf{3D U-Net}. Once the 3D feature volume is instantiated, we run a 3D U-Net to allow 3D processing. The U-Net reduces the spatial resolution by a factor of 8x before upsampling back to the full 3D resolution.
See a more detailed model diagram in the supplemental.

%\TODO{mention vignette removal?}


% key bits: sample with interpolation; concat the point mask; gravity aligend voxel grid;
% joint slam stream lifting
% multi-stream lifting and compression
% Lifter backbone
% 3d unet 

\subsection{3D Bounding Box Detection}
%
We assume all object bounding boxes are gravity-aligned. We take inspiration from the ImVoxelNet~\cite{rukhovich2022imvoxelnet} to build a proposal-free and anchor-free 3D detection head. 
This head runs on top of the output of \EFMmodel{} and regresses the class distribution ($v^{cls}$) and eight values defining the geometric properties of a 3D bounding box. 
It first predicts a centerness score of $[0,1]$ for each voxel $v$. The centerness ($v^c$) of each voxel defines the probability of containing any object bounding box center.
%
Then the head regresses the seven parameters ($v^{obb}$) that describe the bounding box: three for the height, width and depth of the box, three that predict a small offset for the bounding box center to account for discretization of the voxel grid, and lastly a single yaw rotation parameter that determines the rotation of the bounding box around the gravity direction.
%
The predictions are then filtered with Non-Maximum Suppression (NMS) using 3D Intersection over Union (IoU) and the centerness confidence score. There are three losses used for training, one for each head. The centerness and the classifier heads are trained with focal loss (FL)~\cite{lin2017focal}, and the parameters head is trained using gravity oriented 3D IoU loss. With $N_v$ being the number of voxels:
% We directly supervise with observable 3D bounding boxes for each snippet.
%\TODO{how is the GT centerness defined?} 
% \TODO{Data preparation of GT Obbs in impl.}
% \TODO{shorten no moe than one line; need to explain variables}
\begin{align}\textstyle
    L_{obj} = \frac{1}{N_v}\sum^{N_v}_{n} & w_{c}\text{FL}(v_{n}^{c}, \widehat{v_{n}^{c}}) + w_{iou}\text{IoU}(v_{n}^{obb}, \widehat{v_{n}^{obb}}) + w_{cls}\text{FL}(v_{n}^{cls}, \widehat{v_{n}^{cls}}),
\end{align}


\subsection{3D Surface Regression}
To regress 3D surfaces in the local gravity-aligned voxel grid we regress an occupancy value for every voxel. 
%
We supervise occupancy using the GT depth maps. For each GT depth value we obtain three kinds of occupancy supervisions: one surface point, one free-space point, and one occupied point. We compute the surface point $p_\text{surf}$ by unprojecting the GT depth values using the camera calibrations. The free-space point $p_\text{free}$ is sampled in front of the surface, while the occupied point $p_\text{occ}$ is sampled behind the surface. Both $p_\text{free}$ and $p_\text{occ}$ are randomly sampled up to $\delta$ distance from the surface, which we choose as the length of a voxel.
%
At these locations we sample into the predicted occupancy grid via tri-linear interpolation to obtain predicted values for all three kinds of points. The final loss is then computed via the focal loss~\cite{lin2017focal} by defining the ground-truth probability of free-space point, surface point, and occupied point being $0.0$, $0.5$, and $1.0$, respectively:
\begin{align}
   L_{surf} = \textstyle\frac{1}{N}\sum^{N}_{n} \text{FL}(p_\text{free}^n, 0.0) + \text{FL}(p_\text{surf}^n, 0.5) + \text{FL}(p_\text{occ}^n, 1.0),
\end{align}
where $\text{FL}(\cdot)$ is the focal loss and $N$ is the number of unprojected points from the GT depth maps. Note that we use focal loss to make the optimization to focus less on the easy samples like flat surfaces while more on the hard samples like sharp object boundaries. We further add a total variation (TV) regularization loss on the predicted occupancy volume to encourage local smoothness. 
%to fill in unobserved areas in the volume.

\subsection{Implementation and Training Details}
% We introduce implentation details of~\EFMmodel{} and training details.

\noindent\textbf{Input Snippet Data Preparation}. We sample 1s snippets from the RGB and greyscale video streams of Aria running at 10Hz.
% This yields $T=10$ frames per stream per snippet. 
% We resize the original RGB video to $240\times 240\times3$ and the greyscale video to $320\times240\times1$ which leads to approximately the same angular resolution per pixel. The semi-dense points inputs are causally optimized and only visible within the snippet, which ensures that trained models can operate in an online manner. 
The input resolution for RGB and greyscale videos are $240\times 240\times3$ and $320\times240\times1$, which has approximately the same angular resolution per pixel.
% The points observed in a given frame are \textit{causal}, meaning that the points at time $T_t$ depend only on sensor data from previous timesteps and the current timestep $T_0, T_1, ..., T_T$. This ensures that trained models can operate in an online causal way. 
We leave exploring other image sampling strategies like key-framing~\cite{Xie_2023_ICCV} for future research.
For the GT OBB labels, we use the OBBs which are visible inside the snippet and discretize the centers into the volume to obtain the centerness supervision signal. We further filter out the OBBs which are heavily occluded in the snippet as decribed in Sec.~\ref{sec:dataset}.
% Note that we splat the centerness probability into its neighbours using a Gaussian kernel to facilitate the training convergence. 
%\TODO{confirm the RGB resolution}
%\TODO{add SLAM image details if used}

\noindent\textbf{2D Foundation Model Encoder}. While we could use other 2D foundation models like CLIP~\cite{radford2021learning} to encode the video streams, we find the base DinoV2.5~\cite{oquab2023dinov2,darcet2023vitneedreg} model works well for our tasks and is reasonably efficient. 
%The train with input RGB images are resized to $240\times240\times3$. 
%We use DinoV2.5-base which operates on $14\times14$ patches and outputs features of dimension $17\times17\times768$ for the whole image. In all experiments we freeze the DinoV2.5 weights. The 2D feature maps are upsampled back to the input image $240\times240$ resolution\TODO{check this} using stacked CNN \& upsample layers. See more detailed model diagram in the supplemental.

\noindent\textbf{Feature Volume} We choose a $4m^3$ voxel grid. We adapt the grid resolution to the task. For 3D object detection we find a $6.25cm$ voxel size ($D=H=W=64$) is sufficient and for surface regression we push to $4cm$ ($D=H=W=96$) to be comparable with related work~\cite{sun2021neucon}. 
%The feature dimension of the feature volume is $F$.
%The voxel grid size is chosen as $4\times4\times4$ meter cube, but this is easily configurable. We use a voxel grid resolution of $F\times64\times64\times64$, where $F$ is the final 2D feature dimension, each voxel having a unit length of 6.25 cm on each 3D dimension. \TODO

\noindent\textbf{Hyperparameters}. The 3D Bounding Box Detection head is trained on ASE 10k sequences (with 526k snippets) for $10$ epochs and the 3D Surface Regression head is trained on the ASE 1k sequences (with 54k snippets). Both heads were trained with a base learning rate of $2e^{-4}$ using Adam~\cite{kingma2014adam} optimizer with cosine annealing scheduler. We set $w_{cent}=100$, $w_{iou} = 10$, and $w_{class} = 1.0$ for training 3D bounding box detection. For the inference, we use the centerness threshold of $0.2$ to select the candidate detections and performs simple 3D NMS with radius of $2$ voxels based on the 3D IoU. The focal loss paramters in the 3D Surface training is set to $\alpha=0.25$ and $\gamma=2.0$.

% moved to supplement
% %\subsection{Persisting 3D Predictions}
%We argue that persisting 3D inferences is
%As argued in the introduction, persisting 3D inferences is critical to enable compressed context for AI assitants. 

%Inferring the 3D scene state from individual frame or snippet-level 3D inferences is not only useful for downstream systems but also enables comparing different algorithms and measuring how consistent these 3D inferences are across the full sequence. 
%We describe tracking and fusion systems for both 3D bounding box (3D OBBs) as well as surface predictions to enable accumulation and improvement of 3D OBBs as well as 3D surfaces estimates. 

\subsection{Tracking and Fusion of 3D Bounding Box Predictions}

Sequential object detectors predict an independent set of 3D OBBs per frame (e.g., CubeRCNN~\cite{Brazil_2023_CVPR}) or per snippet (e.g., ImVoxelNet~\cite{rukhovich2022imvoxelnet}). Directly appending all the independent sets of 3D OBBs to the scene state results in noisy and duplicated objects. Therefore, we propose a simple object tracker named \textit{ObbTracker} to track and update the objects in the scene from the sequential predictions.

ObbTracker gets the raw predictions as the inputs and maintains a scene state which consists of a set of tracked 3D OBBs. Note that the raw predictions from the object detector are transformed into the world coordinate system using the camera poses. After filtering out the low-confidence predictions, ObbTracker performs 3 key steps: Association, Update, and Removal, which we will explain in detail.

\noindent \textbf{Association}. We match the predictions that clear a matching score threshold of $p_{assoc}$ and the tracked OBBs in the scene using Hungarian Matching with the cost function:
\begin{align}
C_{assoc} = w_1 \cdot C_{class} + w_2 \cdot C_{bbox2} + w_3 \cdot C_{bbox3} + w_4 \cdot C_{iou2d} + w_5 \cdot C_{iou3d},
\end{align}
where $C_{class}$ measures the discrepancy between the classes of the predictions and the tracked 3D OBBs, which is implemented as the probability of tracked OBB’s class evaluated on the predicted class distribution, $C_{bbox2}$ is the distance between the 2D bounding box centers, $C_{bbox3}$ is the distance between the 3D bounding box centers, $C_{iou2d}$ is the 2D Intersection over Union (IoU) between the 2D bounding boxes, and $C_{ioui3d}$ is the 3D IoU between the 3D bounding boxes. 
Specifically we set $w_1=8$, $w_2=0$, $w_3=1$, $w_4=2$, $w_5=0$ for all experiments.
Then we further filter out the matches where the 2D and 3D IoUs are higher than predefined threshold of $0.2$ for both 2D and 3D IoU.
Those good matches will be used for updating the tracked OBBs. The remaining predictions are added to the scene as new candidates if they clear an instantiation score threshold $p_{inst}$. We tune the instantiation and association score thresholds but typically keep $p_{assoc} = p_{inst} - 0.05$. 
%For \EVL{}, ImVoxelNet, 3DETR we use $p_{inst}=0.3$ and for Cube R-CNN $p_{inst} = 0.55$.

\noindent \textbf{Update}. The states of the tracked OBBs are updated with the associated predictions using running average. We define the states of an OBB as its scales $s \in \mathbb{R}^3$, pose in the world coordinate frame $T_{obj}^{w} \in SE(3)$, and the class probability distribution $c \in \mathbb{R}^k$ where $k$ is the number of total categories in the taxonomy. 
Given a predicted set of parameters $s'$, $c'$, and ${T_{obj}^{w}}'$, the running average of the object parameters is computed as
\begin{align}
    n' &= n+1 \\
    s &= (s n + s') / n'\\
    c &=  (c n + c') / n'\\
    T_{obj}^{w} &= T_{obj}^{w} \cdot \exp{({T_{obj}^{w}}' \boxminus T_{obj}^{w} / n')} \;,
\end{align}
where ${T_{obj}^{w}}'$ is the predicted object pose, $n$ is the number of observations of the tracked OBB, and $T_a \boxminus T_b = \log(T_a^{-1} T_b)$ is the difference between two poses on the pose manifold SE3 and $\exp{}$ is the exponential map to SE3. 

%Please refer to the supplementary materials for the detailed definition of the $\boxminus$ operator.

\noindent \textbf{Removal}. After the update step, we remove the tracked OBBs which do not get enough observations $w_{min}$ within a certain period of time $t_{inst}$ after it’s added to the scene. We set $n_{min}=2$ and $t_{inst} =1s$.  
We further suppress duplicate scene OBBs by using both 3D non-maximum suppression (NMS) with IoU threshold of $0.1$ and 2D NMS with IoU threshold of $0.5$.

\subsection{Fusion of Surface Predictions}
To persist local surface observations into a globally consistent estimation, a fusion system is needed for either depth-based or volumetric methods.

\noindent \textbf{Depth map fusion.} When evaluating depth-map-based surface reconstruction, we use a simple implementation from~\cite{zeng20163dmatch} to fuse per-frame observations into an TSDF volume~\cite{newcombe2011kinectfusion}. The method iteratively projects global volume to depth maps, and computes TSDF values by merging the current depth observations in a global averaging way. After getting the global TSDF volume, the meshes are extracted using the marching cubes algorithm~\cite{lorensen1998marching}. 

\noindent \textbf{Occupancy fusion.} For volumetric methods like \EFMmodel{}, a local occupancy volume is predicted per snippet. We fuse the local occupancy grids into a global grid. Each local occupancy grid is iteratively fused into the global volume, by transforming the global volume to local ones and sampling from them. Similar to depth map fusion, We extract mesh surfaces using marching cubes algorithm.

\noindent \textbf{Implementation details.} Both surface fusion algorithms follow an iterative style by feeding the current observations to a global TSDF (for depths) or occupancy (for \EFMmodel{}) volume. We use the iso-level of $0$ for TSDF fusion and $0.5$ for occupancy fusion. In surface extraction, we set a minimal number of observations for a voxel to be considered valid, which is $2$ for TSDF fusion and $5$ for occupancy fusion.

% For the volumetric occupancy or TSDF predictions we resample local grid predictions into a global scene-wide grid and compute running averages. After this fusion we extract meshes using marching cubes~\cite{lorensen1998marching} with a iso-level of $0.5$ for occupancy and $0$ for TSDF fusion. 

